\subsection{Reduced Norms and Traces}

Recall that we denote by A a central simple algebra over the commutative field K of reduced degree n.

PROPOSITION 2. --- Let a be an element of A and \(\mathrm{Pc}(a;\mathrm{X})\) its characteristic polynomial. There exists a unique monic polynomial P in K{[}X{]} such that we have \(\mathrm{Pc}(a;\mathrm{X})=\mathrm{P}(\mathrm{X})^{n}\) .

\begin{enumerate}
\def\labelenumi{\Alph{enumi})}
\tightlist
\item
  The uniqueness of P follows from Lemma 2 below.
\end{enumerate}

Lemma 2. --- Let P and Q be monic polynomials in K{[}X{]} and s is a strictly positive integer. If \(P^{s} = Q^{s}\) , then we have P = Q.

Let \(\mathcal{I}\) be the set of irreducible monic polynomials in K{[}X{]}. Since P and Q are monic, there exist elements \((a_{\mathrm{F}})\) and \((b_{\mathrm{F}})\) of \(\mathbf{N}^{(\mathcal{I})}\) such that we have \(\mathrm{P} = \prod_{\mathrm{F} \in \mathcal{I}} \mathrm{F}^{a_{\mathrm{F}}}\) and \(\mathrm{Q} = \prod_{\mathrm{F} \in \mathcal{I}} \mathrm{F}^{b_{\mathrm{F}}}\) . It follows from the equality \(\mathrm{P}^s = \mathrm{Q}^s\) and the uniqueness of the decomposition into irreducible factors that we have \(sa_{\mathrm{F}} = sb_{\mathrm{F}}\) for every \(\mathrm{F} \in \mathcal{I}\) . Since \(s\) is strictly positive, we consequently have \(a_{\mathrm{F}} = b_{\mathrm{F}}\) for every \(\mathrm{F} \in \mathcal{I}\) , and therefore \(\mathrm{P} = \mathrm{Q}\) .

\begin{enumerate}
\def\labelenumi{\Alph{enumi})}
\setcounter{enumi}{1}
\tightlist
\item
  Let us now prove the existence of P.
\end{enumerate}

By Theorem 1 of VIII, p.~252, there exist a Galois extension L of K of finite degree and an isomorphism of L-algebras \(\theta : \mathrm{A}_{(\mathrm{L})} \to \mathbf{M}_n(\mathrm{L})\) . By formula (12) of III, §9, No.~1, p.~542, the polynomial \(\mathrm{Pc}(a; \mathrm{X})\) is also the characteristic polynomial of the element \(1 \otimes a\) of the L-algebra \(\mathrm{A}_{(\mathrm{L})}\) , hence of the element \(\theta(1 \otimes a)\) of the L-algebra \(\mathbf{M}_n(\mathrm{L})\) .

Set \(\mathrm{P}(\mathrm{X})=\det(\mathrm{XI}_{n}-\theta(1\otimes a))\) ; it is a monic polynomial in L{[}X{]}. By Example 3 of III, §9, No.~3, p.~545, we have

\[
\mathrm{Pc} (a; \mathrm{X}) = \mathrm{P} (\mathrm{X}) ^ {n}.\tag{19}
\]

Let G be the Galois group of L over K. For \(\sigma\in G\) , denote by \(\overline{\sigma}\) the automorphism of the ring L{[}X{]} that coincides with \(\sigma\) on L and fixes X. Then K{[}X{]} is the set of polynomials Q of L{[}X{]} such that we have \(\overline{\sigma}(Q)=Q\) for every \(\sigma\in G\) (V, §10, No.~1, p.~56, Theorem 1). Since the polynomial \(\mathrm{Pc}(a;\mathrm{X})=\mathrm{P}(\mathrm{X})^{n}\) belongs to K{[}X{]}, we have \(\overline{\sigma}(\mathrm{P})^{n}=\mathrm{P}^{n}\) for every \(\sigma\in G\) . By Lemma 2, we therefore have \(\overline{\sigma}(\mathrm{P})=\mathrm{P}\) for every \(\sigma\in G\) , so P belongs to K{[}X{]}.

DEFINITION 1. --- Let a be an element of the algebra A. The reduced characteristic polynomial of a (with respect to A) is the unique monic polynomial in K{[}X{]}, denoted by \(\operatorname{Pcrd}_{\mathrm{A}/\mathrm{K}}(a;\mathrm{X})\) , that satisfies the relation

\[
\mathrm{Pc} _ {\mathrm{A/K}} (a; \mathrm{X}) = \mathrm{Pcrd} _ {\mathrm{A/K}} (a; \mathrm{X}) ^ {n}.\tag{20}
\]

Let a be an element of A. Since A has degree \(n^{2}\) over K, the polynomial \(\mathrm{Pc}_{\mathrm{A/K}}(a;\mathrm{X})\) has degree \(n^{2}\) , so \(\mathrm{Pcrd}_{\mathrm{A/K}}(a;\mathrm{X})\) is a monic polynomial of degree n; let us write it as

\[
\mathrm{Pcrd} _ {\mathrm{A/K}} (a; \mathrm{X}) = \mathrm{X} ^ {n} + \sum_ {r = 0} ^ {n - 1} (- 1) ^ {r} b _ {r} (a) \mathrm{X} ^ {n - r}.\tag{21}
\]

We set

\[
\mathrm{Trd} _ {\mathrm{A/K}} (a) = b _ {1} (a), \quad \mathrm{Nrd} _ {\mathrm{A/K}} (a) = b _ {n} (a).\tag{22}
\]

DEFINITION 2. --- We call \(\operatorname{Trd}_{\mathrm{A/K}}(a)\) the reduced trace of A and \(\operatorname{Nrd}_{\mathrm{A/K}}(a)\) its reduced norm (with respect to the K-algebra A).

When there is no risk of confusion, we leave A and K out of the notation and simply write \(\operatorname{Pcrd}(a; \mathrm{X})\) , \(\operatorname{Trd}(a)\) , and \(\operatorname{Nrd}(a)\) .

The following formulas result from formulas (20) and (22) and formulas (7) and (8) of III, §9, No.~1, p.~542:

\[
\mathrm{Tr} _ {\mathrm{A/K}} (a) = n \mathrm{Trd} _ {\mathrm{A/K}} (a),\tag{23}
\]

\[
\mathrm{N} _ {\mathrm{A/K}} (a) = (\mathrm{Nrd} _ {\mathrm{A/K}} (a)) ^ {n}.\tag{24}
\]

PROPOSITION 3. --- An element a of A is invertible if and only if its reduced norm is nonzero. In particular, A is a field if and only if we have \(\mathrm{Nrd}_{\mathrm{A}/\mathrm{K}}(a) \neq 0\) for every nonzero element a of A.

An element \(a\) of A is invertible if and only if its norm is nonzero (III, §9, No.~4, p.~545, Proposition 3). Proposition 3 therefore follows from formula (24).

Remark. --- Let L be the field \(\mathrm{K}(\mathrm{X})\) of rational fractions in one variable X. The reduced characteristic polynomial of an element a of A is simply the reduced norm of the element \(X \otimes 1 - 1 \otimes a\) of the L-algebra \(\mathrm{A}_{(\mathrm{L})}\) . This follows from the definition of the reduced characteristic polynomial and the formula (III, §9, No.~3, p.~544)

\[
\mathrm{Pc} _ {\mathrm{A/K}} (a; \mathrm{X}) = \mathrm{N} _ {\mathrm{A} _ {(\mathrm{L})} / \mathrm{L}} (\mathrm{X} \otimes 1 - 1 \otimes a).\tag{25}
\]

Examples. --- 1) By Theorem 1 of VIII, p.~120, there exist an integer \(r \geqslant 1\) and a field D such that A is isomorphic to \(\mathbf{M}_{r}(\mathrm{D})\) . Let d be the reduced degree of D over K; we have r = n/d.~Let M be an A-module of finite length \(\ell\) ; we will prove the formula

\[
\mathrm{Pc} _ {M / K} (a_ {M} ;X) = \mathrm{Pcrd} _ {A / K} (a; X) ^ {d\ell}\tag{26}
\]

for every element a of A. The A-module \(A_{s}\) has length r (VIII, p.~121, Lemma 2). The A-modules \(M^{r}\) and \(A_{s}^{\ell}\) have the same length, so they are isomorphic, and we have

\[
\mathrm{Pc} _ {\mathrm{M/K}} (a _ {\mathrm{M}}; \mathrm{X}) ^ {r} = \mathrm{Pc} _ {\mathrm{A/K}} (a; \mathrm{X}) ^ {\ell}
\]

by formula (15) of III, §9, No.~2, p.~542. Since we have rd = n, formula (26) follows from formula (20) and Lemma 2 (VIII, p.~339).

\begin{enumerate}
\def\labelenumi{\arabic{enumi})}
\setcounter{enumi}{1}
\tightlist
\item
  Consider the specific case when A is the algebra \(\operatorname{End}_{\mathrm{K}}(\mathrm{V})\) of endomorphisms of a finite-dimensional K-vector space V. Taking M to be the simple A-module V, we obtain the relations
\end{enumerate}

\[
\begin{array}{c} \mathrm{Pcrd} _ {\mathrm{A/K}} (u; \mathrm{X}) = \chi_ {u} (\mathrm{X}), \\ \mathrm{Nrd} _ {\mathrm{A/K}} (u) = \det (u), \\ \mathrm{Trd} _ {\mathrm{A/K}} (u) = \operatorname{Tr} (u) \end{array}\tag{27}
\]

for every endomorphism u of V.

